% A new book starts from this file. Copy it to <book_dir>/book.tex, run
% setup_book.py and render_cover.py, then build with lualatex (three runs).
% Replace every <<...>>; delete the parts you do not need.
\documentclass[10pt,twoside,openright]{book}
\usepackage{style/animalbook}
\input{style/options.tex}

\begin{document}

\frontmatter
\coverpage{cover/cover.pdf}
\halftitlepage
\copyrightpage{<<draft date, for example 3 November 2026>>}{<<who the book is for;
privacy note if it quotes meetings; cover image credit>>}
\tableofcontents

\chapter{Conventions Used in This Book}

\textbf{Parts.} <<one sentence per part on what it covers, and where a new reader
should start>>

\textbf{Sources.} Every chapter opens with a box saying where its content comes
from. Numbers come only from files in the project: outputs, code, or documents
that cite them. Meetings give context and questions, never numbers.

\textbf{Callouts.}
\keyidea{What to remember from a section.}
\careful{Concepts that are easy to confuse, and common mistakes.}
\workedexample{title}{A number from the project, computed step by step.}
\inproject{Where the concept lives in the repository and the outputs.}
\questionsfor{<<who>>}{What is still open and worth raising.}

\textbf{Evidence tags.} \ev{finding} says something about the question the
project studies; \ev{method check} shows the tools work; \ev{negative} means the
tested effect was not shown; \ev{lead} is a pattern that still needs an
independent test.

\textbf{Priorities.} Sections not yet written are marked P1 (needed now), P2
(useful soon) or P3 (for completeness).

\mainmatter
\part{<<Foundations>>}
\include{chapters/01-<<first-chapter>>}

\part{<<The Problem>>}
\include{chapters/02-<<chapter>>}

\appendix
\include{chapters/<<history-or-glossary>>}

\backmatter
\printindex

\begin{colophon}
The animal on the cover of \emph{<<title>>} is <<the animal, with its Latin
name>>. <<One or two facts about it, and why it fits this book.>>

The cover image is <<artist, date, source and license, or "an engraving
generated with Codex">>.
\end{colophon}

\end{document}
